% =====================================================================
% Page de garde — Cahier des Charges IoT Edge Bridge
% =====================================================================

\thispagestyle{empty}

\begin{titlepage}
  \centering
  \vspace*{1.5cm}

  {\Large \textbf{EPHEC Brussels --- TS6 Technologies de la Santé}}\\[0.4cm]
  {\large Année académique 2025--2026}\\[2.5cm]

  \rule{\linewidth}{0.4pt}\\[0.4cm]
  {\Huge \textbf{Cahier des Charges}}\\[0.4cm]
  {\LARGE Démonstrateur dockerisé}\\[0.4cm]
  {\LARGE \textit{IoT Edge Bridge}}\\[0.2cm]
  {\large Passerelle d'interopérabilité Edge\\
  pour l'intégration de dispositifs médicaux \textit{legacy} au BIHR}\\[0.3cm]
  \rule{\linewidth}{0.4pt}\\[2.0cm]

  \begin{minipage}{0.45\linewidth}
    \begin{flushleft}
      \textbf{Auteur}\\
      Noël Junior Yando Fotso\\
      \texttt{junior.noel@facsciences-uy1.cm}
    \end{flushleft}
  \end{minipage}
  \hfill
  \begin{minipage}{0.45\linewidth}
    \begin{flushright}
      \textbf{Encadrant}\\
      Matthieu Oliviers\\
      \textbf{Date prévisionnelle de remise}\\
      2 mars 2027
    \end{flushright}
  \end{minipage}\\[2.5cm]

  \begin{center}
    \textbf{Version du document : 1.0 --- 4 mai 2026}
  \end{center}

  \vfill

  \footnotesize
  \textit{Ce document est le cahier des charges du démonstrateur logiciel
  associé au projet de fin d'année. Il guide l'implémentation d'une
  passerelle configurable, multi-couches, présentée sous forme d'une
  démonstration entièrement dockerisée. Il ne constitue pas un cahier
  des charges industriel d'un produit commercial.}
\end{titlepage}

% --------------------------------------------------------------------
% Contrôle des versions
% --------------------------------------------------------------------
\chapter*{Contrôle des versions}
\addcontentsline{toc}{chapter}{Contrôle des versions}

\begin{table}[h]
  \centering
  \begin{tabular}{@{}llll@{}}
    \toprule
    \textbf{Version} & \textbf{Date} & \textbf{Auteur} & \textbf{Nature de la révision} \\
    \midrule
    0.1 & 2026-04-20 & N. J. Yando Fotso & Squelette initial, table des matières. \\
    0.2 & 2026-04-28 & N. J. Yando Fotso & Périmètre, vision produit, use cases. \\
    0.9 & 2026-05-02 & N. J. Yando Fotso & Exigences fonctionnelles et non fonctionnelles. \\
    \textbf{1.0} & \textbf{2026-05-04} & \textbf{N. J. Yando Fotso} & \textbf{Première version consolidée pour le jury TS6.} \\
    \bottomrule
  \end{tabular}
  \caption{Journal des versions du présent cahier des charges.}
  \label{tab:cdc-versions}
\end{table}

\noindent\textbf{Statut du document :} version de référence pour le
développement de la démonstration. Toute modification ultérieure est
tracée dans ce tableau et numérotée selon le schéma SemVer
(\textit{majeur.mineur.correctif}).

\noindent\textbf{Distribution :} document interne au projet, partagé
avec le jury et l'encadrant pédagogique.

\clearpage
